\input{../tex-inputs/latex-leseansicht-vorspann}

\section[Arthur Schnitzler an Gustav Schwarzkopf, 2. 9. 1920]{L04406 Arthur Schnitzler an Gustav Schwarzkopf, 2. 9. 1920}
\nopagebreak\mylabel{L04406v}
\rehead{ }\normalsize\beginnumbering\briefempfaengerindex{Schwarzkopf, Gustav@\textsc{Schwarzkopf, Gustav}!zzzSchnitzler, Arthur@\emph{von Arthur Schnitzler}!1920-09-021@{2. 9. 1920}|(be}
\toendnotes[C]{\smallbreak\pagebreak[2]}
\correspDesc{Versand  durch Arthur Schnitzler am 2. 9. 1920 in Altaussee
\newline{}Übermittlung  am 2. 9. 1920 in Altaussee
\newline{}Erhalt  durch Gustav Schwarzkopf im Zeitraum [3. 9. 1920 – 7. 9. 1920?] in Wien}\toendnotes[C]{\smallbreak}
\Standort{DLA, A:Schnitzler, HS.1985.1.1897.}
\physDesc{Bildpostkarte, 413 Zeichen
\newline{}Handschrift: Bleistift, lateinische Kurrent
\newline{}Versand: Stempel: »\nobreak{}\oindex{Altaussee@\textbf{Altaussee}|pwk}Alt Aussee, 2. IX. 20\nobreak{}«.  }\toendnotes[C]{\smallbreak}\pstart{}{\pb}Hr Gustav Schwarzkopf\pend{}\pstart{}Wien I\oindex{I., Innere Stadt@\textbf{I., Innere Stadt}|pw}\pend{}\pstart{}Tiefer Graben 17\oindex{Wien@\textbf{Wien}!I., Innere Stadt@\textbf{I., Innere Stadt}!Tiefer Graben 17@\textbf{Tiefer Graben 17}|pw}.\pend{}{\bigskip}
\pstart
           \centering{}{\pb}\textcolor{gray}{\textbf{Alt-Aussee}}\oindex{Altaussee@\textbf{Altaussee}|pw}\pend
           
\pstart
           \centering{}\textcolor{gray}{\textbf{(709 m), Steiermark\oindex{Steiermark@\textbf{Steiermark}|pw}, mit
                     dem}}\textcolor{gray}{\textbf{Loser}}\oindex{Loser@\textbf{Loser}|pw}\pend
           \vspace{1em}
\pstart
           {\pb}Altaussee, Seevilla\oindex{Seevilla [Altaussee]@\textbf{Seevilla [Altaussee]}|pw}\hfill 2. 9. 1924\pend
           \vspace{0.5em}
\pstart
           lieber Guſtav – seit \label{K_L04406-1v}\edtext{Dinstag auf Reisen}{\lemma{\textnormal{\emph{Dinstag auf Reisen}}}\Cendnote{\textnormal{Vgl. A. S.: \emph{Wiener Schnitzler}, 24. 8. 1920.}}}\label{K_L04406-1}, seit
               \label{K_L04406-2v}\edtext{Freitag, 27.}{\lemma{\textnormal{\emph{Freitag, 27.}}}\Cendnote{\textnormal{Vgl. A. S.: \emph{Wiener Schnitzler}, 27. 8. 1920.}}}\label{K_L04406-2} hier hätte ich heute beinahe meinen ersten regenfreien
               Tag erlebt, we{\geminationn} es nicht eben wieder zu gießen anfinge.
               Trotzd lebt es sich hier leidlich angenehm und, we{\geminationn} man
               durch 30 dividirt spottbillig. Lili\pwindex{Cappellini, Lili 13.\,9.\,1909 Wien – 26.\,7.\,1928 Venedig@\textsc{Cappellini, Lili} (13.\,9.\,1909 Wien – 26.\,7.\,1928 Venedig)|pw} befindet
                  {\pb}sich famos, Heini\pwindex{Schnitzler, Heinrich 9.\,8.\,1902 Hinterbrühl – 12.\,7.\,1982 Wien@\textsc{Schnitzler, Heinrich} (9.\,8.\,1902 Hinterbrühl – 12.\,7.\,1982 Wien)|pw} und Olga\pwindex{Schnitzler, Olga 17.\,1.\,1882 Wien – 13.\,1.\,1970 Lugano@\textsc{Schnitzler, Olga} (17.\,1.\,1882 Wien – 13.\,1.\,1970 Lugano)|pw} sind noch
               ausständig.\pend
           
\pstart
           Beabsichtg \label{K_L04406-3v}\edtext{bis 14. zu bleiben}{\lemma{\textnormal{\emph{bis 14. zu bleiben}}}\Cendnote{\textnormal{Er reiste tatsächlich am 14. 9. 1920 ab, die Heimreise war
                  aber wegen Hochwasser beschwerlich und er kam erst am 16. 9. 1920 in Wien\oindex{Wien@\textbf{Wien}|pwk} an.}}}\label{K_L04406-3}\pend
           \pstart Ihr \spacefill\mbox{Arthur}\pend{}\selectlanguage{ngerman}\endnumbering\briefempfaengerindex{Schwarzkopf, Gustav@\textsc{Schwarzkopf, Gustav}!zzzSchnitzler, Arthur@\emph{von Arthur Schnitzler}!1920-09-021@{2. 9. 1920}|)be}\mylabel{L04406h}
\begin{anhang}
\end{anhang}\newcommand{\dateiname}{L04406}\newcommand{\titel}{Arthur Schnitzler an Gustav Schwarzkopf, 2. 9. 1920}\newcommand{\editorInnen}{Herausgegeben von Jahnke, SelmaMüller, Martin Anton}\input{../tex-inputs/latex-leseansicht-abspann}
